\documentclass[11pt]{exam}
\usepackage{latexsym}
\usepackage{amssymb,amsmath}
\usepackage[pdftex]{graphicx}
\usepackage[top=1in, bottom=1in, left=1in, right=1in]{geometry}
\usepackage{amssymb}
\usepackage{enumitem} 
\usepackage{comment}
\usepackage{mdframed}
%\textwidth=5.5in
%\textheight=8.75in
\pagestyle{headandfoot}
\firstpageheadrule
\runningheadrule
\usepackage{style/header}
\rhead{\name}
\lhead{\class}
\cfoot{Page~\thepage}


\usepackage{latexsym,multirow}
\usepackage{amssymb,amsmath,bm}
\usepackage{graphicx}
\usepackage[color,all,import,arrow]{xy}
\usepackage{booktabs}

%% === bibliography packages ===
\usepackage{natbib}
\bibliographystyle{plainnat}
%% === hyperref options ===
\usepackage{xcolor}
\usepackage[pdftex, bookmarksopen=true, bookmarksnumbered=true,
pdfstartview=FitH, breaklinks=true, urlbordercolor={0 1 0}, citebordercolor={0 0 1}]{hyperref}
\usepackage{colortbl}
\usepackage{subfigure}
\usepackage{subcaption}

%% == tikz
\usepackage{tikz}
\usepackage{inputenc}
\usetikzlibrary{shapes,arrows,trees,fit,positioning,shapes.misc,calc}

% === dcolumn package ===
\usepackage{dcolumn}
\newcolumntype{.}{D{.}{.}{-1}}
\newcolumntype{d}[1]{D{.}{.}{#1}}

% === theorem package ===
\usepackage{theorem}
\theoremstyle{plain}
\theoremheaderfont{\scshape}
\newtheorem{assumption}{Assumption}
\newtheorem{example}{Example}
\newtheorem{definition}{Definition}
\newtheorem{corollary}{Corollary}
\newtheorem{property}{Property}
\newtheorem{proposition}{Proposition}
\newtheorem{theorem}{Theorem}
\newtheorem{remark}{Remark}
\newtheorem{defi}{Definition}
\newtheorem{thm}{Theorem}
\newtheorem{lemma}{Lemma}
\newenvironment{proof}{\paragraph{Proof:}}{\hfill$\square$}

% === some special symbols

\newcommand{\argmax}{\operatornamewithlimits{argmax}}
\newcommand{\argmin}{\operatornamewithlimits{argmin}}
\newcommand{\qed}{\hfill \ensuremath{\Box}}
\newcommand{\ind}{\mbox{$\perp\!\!\!\perp$}}
\newcommand{\nind}{\mbox{$\not\perp\!\!\!\perp$}}
\def\independenT#1#2{\mathrel{\rlap{$#1#2$}\mkern2mu{#1#2}}}
\DeclareMathOperator{\sgn}{sgn}
\DeclareMathOperator{\diag}{diag}
\providecommand{\norm}[1]{\lVert#1\rVert}

\newcommand{\N}{\mathbb{N}}
\newcommand{\E}{\mathbb{E}}
\newcommand{\R}{\mathbb{R}}
\newcommand{\bI}{\bm{I}}
\newcommand{\bH}{\bm{H}}
\newcommand{\bM}{\bm{M}}
\newcommand{\bP}{\bm{P}}
\newcommand{\bQ}{\bm{Q}}
\newcommand{\bV}{\bm{V}}
\newcommand{\bU}{\bm{U}}
\newcommand{\cU}{\mathcal{U}}
\newcommand{\bW}{\bm{W}}
\newcommand{\bw}{\bm{w}}
\newcommand{\bepsilon}{\bm{\epsilon}}
\newcommand{\boldeta}{\bm{\eta}}
\newcommand{\bC}{\bm{C}}
\newcommand{\bl}{\bm{\ell}}
\newcommand{\ATT}{\mathrm{ATT}}
\newcommand{\Yone}{Y(1)}
\newcommand{\Yzero}{Y(0)}
\newcommand{\ATE}{\mathrm{ATE}}
\usepackage{booktabs,tabularx}
\newcommand{\Var}{\mathrm{Var}}
\newcommand{\Cov}{\mathrm{Cov}}

\newcommand{\name}{Section 3: Average Treatment Effects}
\begin{document}
\noindent
\fbox{%
  \begin{minipage}[t]{\linewidth}
    \class, \term \hfill
    \begin{center}
      \textbf{\Large \name}
    \end{center}
    \emph{\tanames} \hfill
  \end{minipage}
}
\begingroup\small

\section{Overview}
We distinguish two complementary perspectives for average treatment effects (ATEs): the finite-population (design-based) view and the super-population (model/sampling-based) view.

\paragraph{Finite-population (design-based) view.}
We fix the $n$ units in our study. By unit-level determinism (SUTVA/consistency assumed), each unit $i$ has fixed potential outcomes $\{Y_i(0),Y_i(1)\}$; the only randomness comes from the assignment mechanism. The causal estimand of interest is the \emph{sample} ATE (SATE),
\[
\tau_{\text{SATE}} \;=\; \frac{1}{n}\sum_{i=1}^n\bigl(Y_i(1)-Y_i(0)\bigr).
\]

\paragraph{Super-population (model-based) view.}
We imagine the units are an i.i.d.\ sample from a large (possibly infinite) population where $(Y(1),Y(0)) \sim \mathbb{P}^\ast$ are random. Randomness comes from the treatment assignment and the sampling. The causal estimand of interest is the \emph{population} ATE (PATE)
\[
\tau_{\text{PATE}} \;=\; \E\!\left[Y(1)-Y(0)\right].
\]

\section{Comparison}

\begin{table}[h!]
\centering
\footnotesize
\begin{tabular}{|>{\centering\arraybackslash}p{0.24\textwidth}|>{\centering\arraybackslash}p{0.35\textwidth}|>{\centering\arraybackslash}p{0.35\textwidth}|}
\hline
 & \textbf{Finite-Population} & \textbf{Super-Population} \\
\hline
\textit{Population} &
Set of units in the study is fixed and finite &
Units are drawn i.i.d.\ from a large or infinite ``super-population.'' \\
\hline
\textit{Potential Outcomes} &
Fixed set $\mathcal{O}_n=\{Y_i(0),Y_i(1)\}_{i=1}^n$ &
$(Y_i(0),Y_i(1)) \text{ i.i.d.\ }\sim \mathbb{P}^{\ast}$ \\
\hline
\textit{Inference} &
Treatment effects for this specific group of units &
Treatment effects for the underlying population distribution \\
\hline
\textit{Source of Randomness} &
Treatment assignment &
\begin{tabular}{@{}c@{}}Treatment assignment \\ Sampling of units from the population
\end{tabular} \\
\hline
\textit{Random Variables} &
Treatment $T_i$ &
Treatment and Potential Outcomes $T_i, Y_i(0), Y_i(1)$ \\
\hline
\textit{Example} &
What is the ATE on these $n$ units in our experiment? &
What is the ATE in the population from which these $n$ units were sampled? \\
\hline
\textit{Validity} &
Internal &
External \\
\hline
\textit{Estimand} &
$\displaystyle \tau_{\mathrm{SATE}}=\frac{1}{n}\sum_{i=1}^n\!\big(Y_i(1)-Y_i(0)\big)$ &
$\displaystyle \tau_{\mathrm{PATE}}=\mathbb{E}\!\left[Y(1)-Y(0)\right]$ \\
\hline
\textit{Estimator} & \multicolumn{2}{c|}{%
\begin{tabular}{c}
Difference in means \\
$\hat\tau = \overline Y_1 - \overline Y_0$
\end{tabular}} \\
\hline
\textit{Unbiasedness} &
$\E[\hat\tau \mid \mathcal{O}_n] = \tau_{\mathrm{SATE}}$ & $\E[\hat\tau] = \tau_{\mathrm{PATE}}$
\\
\hline
% main variance row
\textit{Variance of Estimator} &
$\displaystyle \Var(\hat\tau\mid \mathcal O_n)
= \frac{1}{n}\!\left(\frac{n_0}{n_1}S_1^2+\frac{n_1}{n_0}S_0^2+2S_{01}\right)$
&
$\displaystyle \Var(\hat\tau)=\frac{\sigma_1^2}{n_1}+\frac{\sigma_0^2}{n_0}$ \\
\hline
% bounds row (merged across the two framework columns)
\textit{Estimator of Variance} &
\multicolumn{2}{c|}{%
\begin{tabular}{c}
$\displaystyle \widehat{\Var}(\hat\tau)=\frac{\hat\sigma_1^2}{n_1}+\frac{\hat\sigma_0^2}{n_0}$\\
Finite-Population: Conservative estimator $\mathbb{E}[\widehat{\Var}(\hat\tau) \mid \mathcal O_n] \geq \Var(\hat\tau \mid \mathcal O_n)$, 
\\
equal to design variance iff $Y_i(1)-Y_i(0)=c$ for all $i$, \\
Super-Population: Unbiased estimator $\ \mathbb{E}[\widehat{\Var}(\hat\tau)]=\Var(\hat\tau)$
\end{tabular}} \\
\hline
% Conservative:\ $\displaystyle \widehat{\Var}(\hat\tau)=\frac{\hat\sigma_1^2}{n_1}+\frac{\hat\sigma_0^2}{n_0}$ \ (equal to design variance iff $Y_i(1)-Y_i(0)=c$ for all $i$) &
% Unbiased:\ $\displaystyle \widehat{\Var}(\hat\tau)=\frac{\hat\sigma_1^2}{n_1}+\frac{\hat\sigma_0^2}{n_0}$, and $\ \mathbb{E}\{\widehat{\Var}(\hat\tau)\}=\Var(\hat\tau)$ \\
% \hline
\textit{Asymptotics \qquad \tiny{$\frac{n_1}{n} \to p \in (0, 1)$}} &
\begin{tabular}{@{}c@{}}
Finite-population CLT:\\
$\displaystyle 
\frac{\hat\tau - \tau_{\mathrm{SATE}}}
{\sqrt{\tfrac{1}{n}\!\left(\tfrac{n_0}{n_1}S_1^2+\tfrac{n_1}{n_0}S_0^2+2S_{01}\right)}}
\;\xrightarrow{d}\; \mathcal{N}(0,1)$
\end{tabular}
&
\begin{tabular}{@{}c@{}}
Super-population LLN and CLT:\\
Consistency (LLN): $\hat\tau \xrightarrow{p} \tau_{\mathrm{PATE}}$\\
CLT: $\displaystyle 
\frac{\hat\tau - \tau_{\mathrm{PATE}}}
{\sqrt{\tfrac{\sigma_1^2}{n_1}+\tfrac{\sigma_0^2}{n_0}}}
\;\xrightarrow{d}\; \mathcal{N}(0,1)$
\end{tabular} \\
\hline
\textit{CI in practice} & \multicolumn{2}{c|}{%
\begin{tabular}{c}
$\displaystyle \text{s.e.}=\sqrt{\frac{\hat\sigma_1^2}{n_1}+\frac{\hat\sigma_0^2}{n_0}}$\\
$\displaystyle (1-\alpha)\times 100\%\ \text{CI: }\big[\hat\tau \pm z_{1-\alpha/2}\cdot \text{s.e.}\big]$
\\
Conservative in finite-population; Unbiased in super-population
\end{tabular}} \\
\hline
\end{tabular}

\vspace{0.5em}
\noindent\small\textbf{Assumptions}
\\
The Table is valid under SUTVA and a completely randomized design (CRD).

\noindent\small\textbf{Notation}
\[
\begin{aligned}
&n_1=\sum_{i=1}^n T_i,\quad n_0=n-n_1\\
&\overline{Y(t)} =\frac{1}{n}\sum_{i=1}^n Y_i(t), \qquad \overline Y_t=\frac{1}{n_t}\sum_{i:T_i=t}Y_i, \\
&S_t^2=\frac{1}{n-1}\sum_{i=1}^n\big(Y_i(t)-\overline{Y(t)}\big)^2,\quad 
S_{01}=\frac{1}{n-1}\sum_{i=1}^n\!\big(Y_i(1)-\overline{Y(1)}\big)\big(Y_i(0)-\overline{Y(0)}\big),\\
&\sigma_t^2=\Var\!\big(Y_i(t)\big) = \E[S_t^2],\qquad \hat\sigma_t^2=\frac{1}{n_t-1}\sum_{i:T_i=t}\big(Y_i-\overline Y_t\big)^2.\\
\end{aligned}
\]

\noindent\small\textbf{Finite-Population Bounds for the Variance using Cauchy Schwarz}

Since $S_{01}$ is unobserved, we can obtain these (tight without further assumptions) bounds on the variance for the $\tau_{\mathrm{SATE}}$
\[
\displaystyle \frac{n_0n_1}{n}\!\left(\frac{S_1}{n_1}-\frac{S_0}{n_0}\right)^{\!2}
\le \Var(\hat\tau\mid \mathcal O_n) \le
\frac{n_0n_1}{n}\!\left(\frac{S_1}{n_1}+\frac{S_0}{n_0}\right)^{\!2}
\]
\end{table}

\newpage
\section{Results}
\begin{lemma}[Unbiasedness of the difference-in-means estimator]
\label{lem:unbiasedness}
Let $\hat\tau=\overline Y_1-\overline Y_0$, where $\overline Y_t=\frac{1}{n_t}\sum_{i:T_i=t}Y_i$ and $n_t=\sum_{i=1}^n\mathbf 1\{T_i=t\}$.
\begin{enumerate}[label=(\roman*)]
    \item (\textbf{Finite population / SATE}) Under a completely randomized design (CDR) with fixed $n_1$ and $n_0$, conditioning on the fixed potential outcomes $\mathcal O_n=\{Y_i(0),Y_i(1)\}_{i=1}^n$,
    \[
    \mathbb{E}\!\left[\hat\tau\mid \mathcal O_n\right] \;=\; \tau_{\mathrm{SATE}}
    \;=\; \frac{1}{n}\sum_{i=1}^n\big(Y_i(1)-Y_i(0)\big).
    \]
    \item (\textbf{Super-population / PATE}) If, in addition, the $n$ units are drawn as an i.i.d.\ simple random sample from a population where $(Y(1),Y(0))$ are random, then
    \[
    \mathbb{E}\!\left[\hat\tau\right] \;=\; \tau_{\mathrm{PATE}}
    \;=\; \mathbb{E}\!\big[Y(1)-Y(0)\big].
    \]
\end{enumerate}
\end{lemma}

\begin{remark}
    We also have that $\E[\hat\sigma_t^2] = \sigma_t^2$ by the usual unbiasdness of the sample variance.
\end{remark}

\begin{proof}
(i) \emph{Finite population.} Using linearity of expectation and the fact that the potential outcomes are fixed given $\mathcal O_n$,
\[
\mathbb{E}\!\left[\overline Y_1\mid \mathcal O_n\right]
= \mathbb{E}\!\left[\frac{1}{n_1}\sum_{i=1}^n T_i\,Y_i(1)\,\Big|\, \mathcal O_n\right]
= \frac{1}{n_1}\sum_{i=1}^n \,Y_i(1) \mathbb{E}\!\left[ T_i\,\mid\, \mathcal O_n\right]
= \frac{1}{n_1}\sum_{i=1}^n \,Y_i(1) \frac{n_1}{n}
= \frac{1}{n}\sum_{i=1}^n Y_i(1).
\]
The equality follows because, by symmetry of the CRD with fixed $n_1$, each $i$ is selected into treatment with probability $n_1/n$, 
\[
\mathbb{E}\!\left[ T_i\, \mid \, \mathcal O_n\right] = \mathbb{E}\!\left[ T_i\right] = \Pr(T_1 = 1) = \frac{n_1}{n}
\]
Analogously,
\[
\mathbb{E}\!\left[\overline Y_0\mid \mathcal O_n\right]
= \frac{1}{n}\sum_{i=1}^n Y_i(0).
\]
Subtracting yields
\[
\mathbb{E}\!\left[\hat\tau\mid \mathcal O_n\right]
= \frac{1}{n}\sum_{i=1}^n \big(Y_i(1)-Y_i(0)\big)
= \tau_{\mathrm{SATE}}.
\]

(ii) \emph{Super-population.} Using the law of iterated expectation and (i),
\[
\mathbb{E}[\hat\tau]
= \mathbb{E}\!\left[\;\mathbb{E}[\hat\tau \mid \mathcal O_n]\;\right]
= \mathbb{E}\!\left[\frac{1}{n}\sum_{i=1}^n \big(Y_i(1)-Y_i(0)\big)\right]
= \mathbb{E}\!\big[Y(1)-Y(0)\big]
= \tau_{\mathrm{PATE}},
\]
where the last equality uses i.i.d.\ sampling so that the expectation of the sample mean equals the population mean.
\end{proof}

\begin{lemma}[Variance of the difference-in-means estimator]
\label{lem:var}
Let $\hat\tau=\overline Y_1-\overline Y_0$, where $\overline Y_t=\frac{1}{n_t}\sum_{i:T_i=t}Y_i$ and $n_t=\sum_{i=1}^n\mathbf 1\{T_i=t\}$, under a completely randomized design (CRD) with fixed $n_1$ and $n_0$.
\begin{enumerate}[label=(\roman*)]
    \item (\textbf{Finite population / SATE}) Conditioning on $\mathcal O_n=\{Y_i(0),Y_i(1)\}_{i=1}^n$,
    \[
    \Var(\hat\tau \mid \mathcal O_n)
    \;=\; \frac{1}{n}\!\left(\frac{n_0}{n_1}S_1^2+\frac{n_1}{n_0}S_0^2+2S_{01}\right) = \frac{S_1^2}{n_1}+\frac{S_0^2}{n_0}-\frac{S_\ast^2}{n} \leq \frac{S_1^2}{n_1}+\frac{S_0^2}{n_0},
    \]
    with $S_\ast^2=S_1^2+S_0^2-2S_{01} \geq 0$ and $\E\left[S_\ast^2\right] = \Var\!\big(Y_i(1)-Y_i(0)\big)$.
    \item (\textbf{Super-population / PATE}) If, in addition, the $n$ units are drawn i.i.d.\ from a population where $(Y(1),Y(0))$ are random,
    \[
    \Var(\hat\tau)
    \;=\; \frac{\sigma_1^2}{n_1}+\frac{\sigma_0^2}{n_0}\,.
    \]
\end{enumerate}
\end{lemma}

\begin{remark}
    The above lemma illustrates why we can use the estimator for the super-population variance for the finite-population variance. In expectation this estimator yields.
    \[
    \mathbb{E}[\widehat{\Var}(\hat\tau) \mid \mathcal O_n] = \mathbb{E}\left[\frac{\hat\sigma_1^2}{n_1}+\frac{\hat\sigma_0^2}{n_0}\right] = \frac{S_1^2}{n_1}+\frac{S_0^2}{n_0} \geq \Var(\hat\tau \mid \mathcal O_n)
    \]
\end{remark}

\begin{proof}
(i) \emph{Finite population.} 
The first equality follows from the review questions. We show the second equality here. By definition $2S_{01} = S_1^2 + S_0^2 - S_\ast^2$. Inserting into the variance
\begin{align*}
     \Var(\hat\tau \mid \mathcal O_n)
    &= \frac{1}{n}\!\left(\frac{n_0}{n_1}S_1^2+\frac{n_1}{n_0}S_0^2+2S_{01}\right)
    \\
    &= \frac{1}{n}\!\left(\frac{n_0}{n_1}S_1^2+\frac{n_1}{n_0}S_0^2+S_1^2 + S_0^2 - S_\ast^2\right)
    \\
    &= \frac{1}{n}\!\left(\frac{n_0 + n_1}{n_1}S_1^2+\frac{n_1 + n_0}{n_0}S_0^2 - S_\ast^2\right)
    \\ 
    &= \frac{1}{n}\!\left(\frac{n}{n_1}S_1^2+\frac{n}{n_0}S_0^2 - S_\ast^2\right)
    \\
    &= \frac{S_1^2}{n_1}+\frac{S_0^2}{n_0}-\frac{S_\ast^2}{n}. 
\end{align*}
This proves the second equality. It remains to show that $S_\ast^2 \geq 0$. By Cauchy-Schwarz we have
\[
S_{01} \leq S_1 S_0.
\]
Noting that $S_t \geq 0$ we have 

\[S_\ast^2=S_1^2+S_0^2-2S_{01} \geq S_1^2+S_0^2 - 2S_1S_0 = (S_1 - S_0)^2 \geq 0,
\]
which proves the last claim.


(ii) \emph{Super-population.} By the law of iterated iterated variance and our i.i.d. assumption:
\begin{align*}
\Var(\hat\tau)&=\E\!\big[\Var(\hat\tau\mid \mathcal O_n)\big]+\Var\!\big(\E[\hat\tau\mid \mathcal O_n]\big)
\\
&= \E \left[\frac{S_1^2}{n_1}+\frac{S_0^2}{n_0}-\frac{S_\star^2}{n}\right] + \Var\left(\frac{1}{n}\sum_{i=1}^n\!\big(Y_i(1)-Y_i(0)\big)\right)
\\
&= \frac{\sigma_1^2}{n_1}+\frac{\sigma_0^2}{n_0} - \frac{1}{n}E[S_\star^2] + \frac{1}{n} \Var(Y_i(1)-Y_i(0))
\end{align*}
since $\E[S_t^2]=\sigma_t^2$ by the unbiasedness of the sample variance. It remains to show that $E[S_\star^2] = \Var(Y(1)-Y(0))$.
Let $\mu_1=\E[Y(1)]$, $\mu_0=\E[Y(0)]$. Expand
\begin{align*}
    S_{01} = \frac{1}{n-1} \sum_{i=1}^n\big(Y_i(1)-\overline{Y(1)}\big)\big(Y_i(0)-\overline{Y(0)}\big)
=\frac{1}{n-1} \left[\sum_{i=1}^n\big(Y_i(1)-\mu_1\big)\big(Y_i(0)-\mu_0\big)
- n\big(\overline{Y(1)}-\mu_1\big)\big(\overline{Y(0)}-\mu_0\big)\right]
\end{align*}
Taking expectations,
\begin{align*}
    \E\!\left[\sum_{i=1}^n\big(Y_i(1)-\mu_1\big)\big(Y_i(0)-\mu_0\big)\right]
&= n \Cov(Y_i(1), Y_i(0))
\\
\E\!\left[n\big(\overline{Y(1)}-\mu_1\big)\big(\overline{Y(0)}-\mu_0\big)\right] &= n\,\Cov\!\big(\overline{Y(1)},\overline{Y(0)}\big)
\\
&=n \frac{1}{n^2}\sum_{i, j =1}^n \Cov(Y_i(1), Y_j(0))
\\
&= \frac{1}{n} \sum_{i=1}^n \Cov(Y_i(1), Y_i(0))
\\
&= \Cov(Y_i(1), Y_i(0))
\end{align*}
Therefore
\[
\E\!\left[S_{01}\right]=\frac{1}{n-1} (n-1) \Cov(Y_i(1), Y_i(0)) = \Cov(Y_i(1), Y_i(0))
\]
By linearity of expectation,
\[
\E[S_\ast^2]
= \E[S_1^2]+\E[S_0^2]-2\,\E[S_{01}]
= \sigma_1^2+\sigma_0^2-2\Cov(Y_i(1), Y_i(0)).
\]
Finally, for any two random variables $A,B$,
\(\Var(A-B)=\Var(A)+\Var(B)-2\Cov(A,B)\).
With $A=Y(1)$ and $B=Y(0)$,
\[
\Var\!\big(Y(1)-Y(0)\big)=\sigma_1^2+\sigma_0^2-2\Cov(Y_i(1), Y_i(0)).
\]
Thus $\E[S_\ast^2]=\Var\!\big(Y(1)-Y(0)\big)$, as claimed.




\end{proof}







\endgroup
\end{document}
